
%==============================================
\title{Concentration Effect in Youth Crime}
%==============================================

\section{Research question}

\noindent Does the concentration of young people leaving school at the same time
affect crime in the hours after release?

\paragraph{Reduced form.} Did the 2014 change in the schedule of 1st-year EM
classes, which stopped them from leaving school at the same time as the 2nd
and 3rd years, change crime by or against 2nd- and 3rd-year students after
release, in schools that used to release all EM years together?

 \section{Motivation}

School reduces youth crime in two ways: it raises human capital in the long
run \citep{lochner2004}, and it keeps students in a supervised environment, where they can not commit crime 
(incapacitation). But school also brings young people
together, and that can increase crime. \citet{billings2019} use the redrawing of
school boundaries in Charlotte-Mecklenburg in 2002 and find that one standard
deviation more same-school, same-grade, same-race and same-gender peers living
within 1 km raises the probability of arrest at ages 16--21 by 3.9 percentage
points (23\%). Two students who live within 1 km of each other are about six
times more likely to be arrested \emph{together} if they are assigned to the
same school.

Concentration can also raise crime in the short run, while students are
together. The literature shows evidence of the concentration effects through variation in session-days and the
type of crime. \citet{jacob2003} show that propery crime falls by 14\% on school days, but juvenile
violent crime rises by 28\%, and the violent crime peaks at about 3 p.m.,
when school lets out. The authors arugue that the fall in property crime is an evidence of incapacitation effects,
 and the rise of violent crime is an effect of concentration.  \citet{luallen2006} follows the same line and shows that the
property crime increases
$+29\%$ and violent crime $-32\%$ on strike days. 


This is weak evidence of a concentration effect for some reasons.
First, they identify only the net effect of incapacitation and concentration.
The effect of school on violent crime is the sum of incapacitation ($-$) and
concentration ($+$). Since the net effect is positive (+28\%), concentration
outweighs incapacitation for violent crime. The effect on property crime is the
sum of incapacitation ($-$) and concentration, whose sign is unknown. If we
assume it is non-negative, the negative net effect ($-14\%$) means that
incapacitation outweighs concentration for property crime. Under these sign
assumptions, they identify only the lower bound of the concentration effect. 


Second, the source of variation is a whole day without school, which changes
substantially what students do during the day, not only how many other students
they meet. The changes in property and violent crime can therefore reflect the
different opportunities students face on an ``atypical'' day: the same student
can change which crime they commit. For example, a student who would have
fought at school (violent crime) may instead shoplift at a mall (property
crime), because a free day leaves time to comute to the mall that is far away.
 If students substitute across crimes in this way, the type of
crime does not reveal whether incapacitation or concentration is at work. This
fits the results of \citet{jacob2003}: on in-service days, total serious
juvenile crime is almost unchanged, because the rise in property crime and the
fall in violent crime roughly cancel out.

Third (and less severe), although \citet{jacob2003} rules out the bias of school-reporting by 
estimating the effect of non school hours crimes or crimes that occured out of schools, their effect decrease with their significance. 

Separating concentration from incapacitation is hard. Any policy that changes
when students leave school usually also changes how long they stay there, and
so how much time they spend supervised. This paper uses a source of variation
that changes concentration while keeping school and supervision fixed: whether
9th graders are released at the same time as high-school students in their own
school. We compare classes that have school on the same day and spend the same
number of hours in class. Only the timing of release differs, and with it the
number and the age of the students who leave with them. This design has four
advantages over the existing evidence:
\begin{enumerate}
  \item \textbf{Incapacitation is held fixed.} School is in session on both
        sides of the comparison, and class duration is the same. Treated and
        control students have the same hours of supervision and the same free
        time. Any difference in crime comes from who they are released with,
        not from how long they were in school.
  \item \textbf{Concentration is measured, not inferred.} We observe how many
        students leave the school at the same time, and within any window
        around the bell. This gives a dose--response: crime as a function of
        the size of the crowd at the gate, not just the sign of an effect.
  \item \textbf{Size and composition can be separated.} Sharing release time
        with high-school students raises both the number of youths on the
        street and their age. Comparing exposure to older students with
        exposure to the same number of same-age students separates a pure
        concentration effect from a peer effect of older students.
  \item \textbf{It targets the right margin.} The effect is measured on the
        street after release, outside school grounds, where reporting by school
        staff plays no role. That is the margin a change in bell times acts on.
\end{enumerate}
The setting adds external value. Most of the evidence comes from US cities. In
Brazil, public schools run several shifts per day, so release times vary a lot
within and across schools. Youth violence is also much higher, so the stakes of
getting the policy right are larger.



Understanding the contemporaneous effect of concentration matters for two reasons:
\begin{enumerate}
  \item It is cheap to prevent. If shared release raises crime, staggering bell
        times is a low-cost policy, unlike lengthening the school day or
        reassigning students across schools.
  \item It may feed the long-run effect. If crimes committed together after
        school are where partnerships start, the short-run concentration effect
        is one channel behind the networks documented by \citet{billings2019}. \citet{luallen2006} finds that juveniles who become repeat
        offenders are more likely to commit their first crime on a strike day.
\end{enumerate}

We also investigate:
\begin{enumerate}
  \item \textbf{Displacement.} Does staggered release reduce crime, or only
        shift it to other times and days? \citet{jacob2003} and \citet{luallen2006} find no temporal displacement of violent crime (leads and lags,
        and weekly and monthly totals). We can run the same tests.
  \item \textbf{Place.} Where does concentration-driven crime happen: at the
        school gate, or along the routes between school and home?
  \item \textbf{Type of crime.} Is the effect driven by fights between students
        (assaults with young victims) or by property crime? Is it larger for
        crimes with several juvenile offenders? \citet{jacob2003} find no
        evidence that school helps juveniles coordinate crimes. \citet{billings2019} find that partnerships are common among older
        youth.
      \item Crime against students
\end{enumerate}

 \section{Mechanisms}

I can think of three mechanisms on which conecetration of youth can increase crime. First, is the numeric concentration of people, second the peer effect and 
third the report. Each of these mechanisms operate in the different manner. 
\begin{enumerate}
      \item[A.] First, the number of people: how many young people are around.
            A concentration of people can increase crime because:
      \begin{enumerate}
            \item there are more interactions, which can help youths
                  coordinate crimes. \citet{jacob2003} find no evidence that
                  school increases the coordination of crime: crimes with
                  several offenders do not respond more to school days. We can
                  test whether the number of crimes committed by more than one
                  student increases.
            \item there are more people who can fight each other. We can test
                  this, but we cannot separate it from old rivalries (the peer
                  effect below).
            \item there is less adult supervision per youth (police, school
                  staff, security guards on the street), as long as the police
                  do not respond to the crowd.
            \item more students can attract offenders, which increases crime
                  against students. We can check whether the number of student
                  victims increases.
      \end{enumerate}
      \item[B.] Second, the peer effect: .
      \begin{enumerate}
            \item \textbf{Learning.} Students can learn from other students.
                  In our main specification, we are measuring the exposition of older students to
                  younger ones so this channel may be weaker in our case. But we can test (long-shot) if there is an increase in crimes
                   committed jointly with a prior offernder and a first-time offender.  Perhaps we do not have enough variation. 
            \item \textbf{Recruitment.} If a 3rd-year
                  student pulls 1st-year students into a theft, the 3rd-year
                  student's own crime also increases, because they offend
                  together. We can test this with crimes committed by mixed-grade
                  groups.             
            \item \textbf{Status.} Older students may act out/show off to
                  impress younger ones. 
            \item \textbf{Within student rivalries.} Students can fight because of their
                  rivalries, including hazing of 1st-year students by older
                  students. This predicts fights
                 across grades, with older offenders and younger
                  victims. We can test this. 
      \end{enumerate}
            \item[C.] Third, policing and reporting (I am not sure if we treat this as a mechanism or a confounder... )

            \begin{itemize}
                  \item \textbf{Reporting.} A crowd can change the share of
                        crimes that reach the police, in either direction.
                        It can decrease reporting:                       offenders are harder to identify in a crowd and
                         each witness expects
                        someone else to call
                   Also, fights between students that would otherwise happen away from the school
may happen at the school gate when students are released together. There, the
school is more likely to handle them as a discipline matter, so they are less
likely to be reported to the police.
                         It can increase reporting: there are
                        more witnesses and more phones to film or call.
                  \item \textbf{More police.} When students leave together,
                        the police can respond with more patrols at the school
                        gate. This has two effects: more
                        police can deter crime (a real change in
                        behavior), and more police can detect more, 
            \end{itemize}
            We can test for these channels by checking whether, with shared
            release:
            \begin{enumerate}
                  \item the average time between the crime and the arrival of
                        the police (or the report) changes;
                  \item the share of crimes in which the police are already at
                        the scene, or arrest the offender in the act
                        (\emph{flagrante}), increases;
                  \item the share of incidents started by the police (patrol)
                        rather than by a call from a victim or witness changes;
                  \item serious crimes, which are reported whoever is watching
                       , respond in the same way as
                        minor crimes. If only minor crimes respond, the effect
                        is more likely to come from reporting.
            \end{enumerate}
\end{enumerate}




\section{Population of Study}

\begin{itemize}
  \item \textbf{Population.} Students in the 2nd and 3rd years of high school
        (EM) in state schools in Minas Gerais that offer all three years of EM.
  \item \textbf{Treatment.} The class is released at the same time as the
        1st-year EM classes of the same school and shift. A continuous version
        is the number of 1st-year students released within $X$ minutes of the
        class.
  \item \textbf{Outcome.} Crime committed by (and against) 2nd- and 3rd-year
        students after release, near the school and in the minutes after the
        bell.
\end{itemize}

\section{Identification Strategy}

\subsection{Motivation for the Identification Strategy}

\paragraph{The naive comparison.} The simplest design compares schools that
release the three EM years together with schools that release each year at a
different time. But release times are not random, and this comparison is
biased by confounders at three levels.

\emph{School level.}
\begin{itemize}
  \item \textbf{Size and density.} Schools with fewer classes can release
        everyone together, so shared release is more common in smaller schools.
        Smaller schools are in less dense areas with less crime. This biases a
        comparison of shared and non-shared schools towards finding that
        concentration \emph{reduces} crime. The bias reverses if treatment is
        measured as the \emph{number} of students released together, which is
        larger in large schools located in dense areas.
  \item \textbf{Infrastructure.} Staggered times often come from lack of rooms
        or extra shifts, which are also related to school quality.
  \item \textbf{Reverse causality.} Principals may stagger release
        \emph{because} of fights at the gate, so treatment responds to past
        violence.
  \item \textbf{Nearby schools.} Other schools releasing at the same time add
        to concentration, and their students may also be affected.
\end{itemize}

\emph{Individual level.}
\begin{itemize}
  \item \textbf{Sorting into schools.} Families choose schools, so students in
        small schools with shared release differ from students in large
        schools (income, distance from home, prior behavior).
  \item \textbf{Sorting into shifts and classes.} Schools assign students to
        shifts and classes by age, repetition, behavior or work status. If
        some classes are released at different times, the students in them may
        differ in their risk of crime.
\end{itemize}

\emph{Time of day.} A different release time is a different clock time, with
different parental supervision, daylight, policing and street activity.

\paragraph{A source of variation: the 2014 change in 1st-year schedules.} In
2014, 1st-year EM classes in state schools in Minas Gerais became about
50 minutes longer, for that year only (probably the sixth period of
\emph{Reinventando o Ensino M\'edio}; to be confirmed). Morning 1st-year classes
moved their release from 11:20--11:30 to about 12:10--12:20, and afternoon
classes from 17:00--17:30 to 17:50--18:20. Start times did not change, and the
schedule of 2nd-year classes did not change either: they kept leaving at
11:20--11:30 and 17:00--17:30. In 2015 the 1st year went back to its previous
schedule.

For a 2nd- or 3rd-year student in a school that released all EM years
together, this means:
\begin{itemize}
  \item \textbf{Concentration falls in 2014 and returns in 2015.} The 1st-year
        students no longer leave with them in 2014, and leave with them again
        in 2015.
  \item \textbf{Their own schedule is unchanged.} Start time, class duration
        and release time are the same, so incapacitation (hours in school,
        free time) and clock time are held fixed \emph{for the treated
        students}.
  \item \textbf{The change is set by the state, not by the school.} It does not
        respond to violence at a particular school's gate.
\end{itemize}

\paragraph{Design.} Compare 2nd- and 3rd-year students in schools that shared
release with the 1st year before 2014 (treated: sharing stops in 2014 and
returns in 2015) with students in schools that did not share release (control:
no change), in 2012--2016. The reversal in 2015 is a second test: if crime
falls in 2014 because of lower concentration, it should rise again in 2015.

\begin{itemize}
  \item \textbf{Comparison:} treated vs control schools, before, during and
        after 2014.
  \item \textbf{Unit:} 2nd- and 3rd-year class (or school $\times$ shift)
        $\times$ year.
  \item \textbf{Fixed effects:} school $\times$ shift and year.
  \item \textbf{Identifying assumption:} without the 2014 change, crime around
        the release of 2nd- and 3rd-year students would have evolved in
        parallel in treated and control schools; and the 2014 change affected
        these students only through the number of students released with
        them.
\end{itemize}

\paragraph{Threats.}
\begin{itemize}
  \item \textbf{Treated and control schools differ.} Schools that shared
        release are smaller (see above). Fixed effects absorb level
        differences, but not different trends. Check pre-trends in 2010--2013.
  \item \textbf{The 1st-year students also change.} In 2014 they leave later
        and spend more time in school, so their own crime changes. The outcome
        must be crime by (or against) 2nd- and 3rd-year students, not all
        youth crime near the school.
  \item \textbf{Other changes in 2014.} The policy may have come with other
        changes at the school (activities, staff) that affect all students. In
        addition, Belo Horizonte hosted World Cup matches in June--July 2014,
        with changes to the school calendar and heavy policing, and 2014 was an
        election year. Year fixed effects absorb what is common to all
        schools, but not effects that differ between treated and control
        schools.
  \item \textbf{The 3rd year is not yet checked.} The unchanged schedule is
        confirmed for the 2nd year (etapas 26/31/36). The 3rd year (etapas
        27/32/37) still needs to be verified.
  \item \textbf{Compliance.} Not all schools may have changed the 1st-year
        schedule in 2014. Use the observed change in each school's schedule,
        not just the year, and count how many schools actually lost shared
        release between the EM years.
\end{itemize}

\paragraph{Descriptive evidence on the 2014 change.} Figure~\ref{fig:tratados}
counts the schools whose main 1st-year morning release time moved 50 minutes or
more later than in the previous year. The change is concentrated in 2014: 170
of 247 schools (69\%). A few schools made the same change in earlier years: 8 from 2011 to 2012 and
11 from 2012 to 2013;
they may have adopted the change before the rest of the network. These 19
schools are dropped from Figures~\ref{fig:horarios_manha}
and~\ref{fig:horarios_tarde}, so the remaining schools are all untreated before
2014.

\begin{figure}[ht]
\centering
\includegraphics[width=0.8\textwidth]{tratados_saida_1em_manha.png}
\caption{Schools whose main 1st-year EM release time in the morning shift moved
50 minutes or more later than in the previous year. Bar height: share of the
sample schools with 1st-year morning classes in both years; number on the bar:
number of schools.}
\label{fig:tratados}
\end{figure}

Figure~\ref{fig:horarios_manha} shows that, in the morning, 1st-year classes
moved from 11:20--11:30 to 12:15--12:20 in 2014 only, while 2nd-year classes
kept leaving at 11:20--11:30. Figure~\ref{fig:horarios_tarde} shows that, in the
afternoon, some schools moved the start time earlier in 2014 (from 13:00 to
12:30). This happens for the 2nd year as well as the 1st, so in the afternoon
the 2014 change does not leave the 2nd-year schedule unchanged in every school.

\begin{figure}[ht]
\centering
\includegraphics[width=\textwidth]{horarios_em2_sem_antecipadas_manha_mapa.png}
\caption{Release time of 1st- and 2nd-year EM classes in the morning shift
(classes starting before 12:00), \% of classes by year. Sample state schools,
excluding the 19 schools that delayed 1st-year release before 2014.}
\label{fig:horarios_manha}
\end{figure}

\begin{figure}[ht]
\centering
\includegraphics[width=\textwidth]{horarios_em2_sem_antecipadas_tarde_entrada_mapa.png}
\caption{Start time of 1st- and 2nd-year EM classes in the afternoon shift
(classes starting between 12:00 and 14:00), \% of classes by year. Sample state
schools, excluding the 19 schools that delayed 1st-year release before 2014.}
\label{fig:horarios_tarde}
\end{figure}

Figure~\ref{fig:hist_diff} compares, within each school, the release time of
the 1st year with that of the 2nd and 3rd years in the morning shift. In 2013
and 2015 almost every school releases all EM years at the same time (difference
of zero). In 2014 most schools release the 1st year about 50 minutes after the
2nd and 3rd years: the median difference is 50 minutes, and 163 of 224 schools
have a difference of 40 minutes or more. Only 37 schools keep releasing all
years together in 2014.

\begin{figure}[ht]
\centering
\includegraphics[width=0.9\textwidth]{hist_diff_saida_1em_manha.png}
\caption{Difference, in each school, between the mean release time of 1st-year
EM classes and the mean release time of 2nd- and 3rd-year EM classes, morning
shift (classes starting before 12:00), in 2013, 2014 and 2015. Positive values:
the 1st year leaves later. Sample state schools with both 1st-year and
2nd- or 3rd-year morning classes in the year. \textbf{Note:} the 19 schools that
delayed 1st-year release by 50 minutes or more before 2014 (8 from 2011 to 2012
and 11 from 2012 to 2013) are excluded. The horizontal axis is cut at $-30$ and $80$ minutes, which
leaves out 3 schools in 2013 (differences of $-120$, $-170$ and $-225$ minutes).}
\label{fig:hist_diff}
\end{figure}




\begin{figure}[ht]
\centering
\resizebox{\textwidth}{!}{%
\begin{tikzpicture}[
    var/.style={rectangle, rounded corners, draw=gray, thick, align=center,
                text width=3.7cm, minimum height=1cm, font=\small,
                execute at begin node={\hyphenpenalty=10000}},
    main/.style={var, draw=black, very thick, fill=gray!10, text width=4cm},
    mech/.style={var, draw=green!50!black, fill=green!5},
    cvar/.style={var, draw=red!75!black, fill=red!4},
    shock/.style={var, draw=blue!70!black, fill=blue!5},
    effect/.style={-{Latex[length=3mm]}, very thick, green!50!black},
    conf/.style={-{Latex[length=2.5mm]}, thick, red!75!black},
    detect/.style={-{Latex[length=2.5mm]}, thick, dashed, orange!85!black},
    iv/.style={-{Latex[length=3mm]}, very thick, blue!70!black},
    threat/.style={-{Latex[length=2.5mm]}, thick, dashed, red!75!black},
    link/.style={thick, red!75!black, {Latex[length=2.5mm]}-{Latex[length=2.5mm]}, dotted}
]
    % ---------------- treatment, outcome, source of variation ----------------
    \node[main]  (conc)  at (0,0)
        {\textbf{Concentration at release}\\1st, 2nd and 3rd years of EM
         released together};
    \node[main]  (crime) at (14,0)
        {\textbf{Crime} by or against 2nd--3rd-year students after release};
    \node[shock] (shock) at (-6,0)
        {\textbf{2014 change}: 1st-year EM classes 50 min longer (one year
         only)};

    % ---------------- mechanisms (concentration) ----------------
    \node[mech] (m1) at (7,6.0) {(1) More encounters, more conflict};
    \node[mech] (m2) at (7,4.4) {(2) Weaker supervision per youth};
    \node[mech] (m3) at (7,2.8) {(3) More targets in one place};
    \node[mech] (m4) at (7,1.2) {(4) Crowd and group behaviour};
    \node[var, draw=orange!85!black, fill=orange!5] (det) at (7,-1.5)
        {(9) More police and witnesses: detection, not behaviour};

    % ---------------- confounders ----------------
    % row 1: direct confounders of concentration and crime
    \node[cvar] (size)   at (-2.5,-4.6) {School size (few classes $\Rightarrow$ shared release)};
    \node[cvar] (infra)  at (1.8,-4.6)  {Infrastructure (rooms, extra shifts)};
    \node[cvar] (past)   at (6.1,-4.6)  {Past violence at the gate (reverse causality)};
    \node[cvar] (sort)   at (10.4,-4.6) {Sorting of students (risk) into schools, shifts, classes};
    \node[cvar] (nearby) at (15.2,-4.6) {Other schools releasing nearby};
    % row 2: deeper causes
    \node[cvar] (dense)  at (1.8,-7.6)  {Density and neighbourhood crime};
    \node[cvar] (clock)  at (10.4,-7.6) {Time of day (parents at work, daylight, policing)};

    % ---------------- source of variation ----------------
    \draw[iv] (shock) -- (conc)
        node[midway, below, font=\scriptsize, text=blue!70!black, align=center]
        {own schedule\\unchanged};
    \draw[threat] (shock.north) .. controls (-6,9) and (14,9) .. (crime.north)
        node[pos=0.5, above, font=\scriptsize, text=red!75!black]
        {1st-years leave later: their own crime changes (keep them out of the outcome)};

    % ---------------- mechanisms ----------------
    \foreach \m in {m1,m2,m3,m4} {
        \draw[effect] (conc) -- (\m.west);
        \draw[effect] (\m.east) -- (crime);
    }
    \draw[detect] (conc) -- (det.west);
    \draw[detect] (det.east) -- (crime)
        node[midway, below, font=\scriptsize, text=orange!85!black] {recorded crime};

    % ---------------- confounder paths ----------------
    \draw[conf] (size) -- (conc);
    \draw[conf] (infra) -- (conc);
    \draw[conf] (infra) -- (crime);
    \draw[conf] (past) -- (conc);
    \draw[conf] (past) -- (crime);
    \draw[conf] (sort) -- (conc);
    \draw[conf] (sort) -- (crime);
    \draw[conf] (nearby) -- (crime);
    % density: drives school size, nearby schools and crime
    \draw[conf] (dense) -- (size);
    \draw[conf] (dense.east) -- (nearby.south west);
    \draw[conf] (dense.south) .. controls (2,-10.2) and (18.5,-10.2) .. (18.5,-3)
        .. controls (18.5,-1) and (17.5,0) .. (crime.east);
    % time of day
    \draw[conf] (clock.east) .. controls (13.5,-7.6) and (14,-3) .. (crime.south);
    \draw[link] (clock.south west) .. controls (6,-10.4) and (-6.5,-9) .. (conc.south west)
        node[pos=0.3, below, font=\scriptsize, text=red!75!black]
        {release time sets clock time};
\end{tikzpicture}%
}
\caption{DAG for the effect of concentration at release on the crime of 2nd-
and 3rd-year EM students. Green: concentration mechanisms (1)--(4), numbered as
in the Mechanisms subsection. Orange dashed: detection (9), which raises
recorded crime without changing behaviour. Red: confounders that open back
doors between concentration and crime (school level, individual sorting, time
of day). Blue: the 2014 change in 1st-year schedules, which moves concentration
for 2nd- and 3rd-year students without changing their own schedule, so it
closes the time-of-day and hours-in-school paths for them; the red dashed arc
is its main threat. The peer mechanisms (5)--(8) are not in this graph: in
this design the students exposed are older than those added to the crowd.}
\label{fig:confounders}
\end{figure}

I have educational data for MG from 2007 to 2017, including proficiency scores


\subsection{Heterogenous Estimation}

Powered effect of concentration: if have two shifts in the smae school. The program makes students go earlier in the school in the afternoon shift (starts at \_)

\subsection{Details:}

\begin{itemize}
      \item Schools with distinct release time of the class within the grade (around 5\%): I would exclude them first, and then think
      \item Windows of the outocme: 1,2 ,3  hours
      \item Effect of morning shift and afternoon shift
\end{itemize}

\subsection{Sources of variations}

\medskip
\noindent\textbf{(a) ``RCT'' within school, class level.}

\noindent\textbf{Level.} Class
 \noindent\textbf{Treatment.} Shared release time with 10th grade
\noindent\textbf{Sample.} Cell = school $\times$ shift $\times$ year
$\times$ class duration. Cells with at least two 9th-grade classes, at
least one released at a time shared with 10th grade and at least one
released at a time not shared with 10th grade. Treated and control
classes have the same duration.

\noindent\textbf{Estimator.} Simple difference in means: 
\begin{equation}
  \hat\beta = \sum_{s} w_s \Big(\overline{V}^{\text{shared}}_{s}
  - \overline{V}^{\text{not shared}}_{s}\Big)
\end{equation}
Equivalent
regression, with $\alpha_s$ a cell fixed effect:
\begin{equation}
  V_{cs} = \alpha_s + \beta\, \text{Shared}_{cs} + \varepsilon_{cs}
\end{equation}

\noindent\textbf{Identifying assumption.} Inside a cell, which class gets
the shared release time is as good as random: students in shared and not
shared classes of the same school, shift and year have the same potential
outcomes. School-level confounders (size, density, neighborhood) and shift
selection are fixed by construction.

\noindent\textbf{Confounders.} 
\begin{itemize}
\item Confounder: the school chooses which class goes to which time
(e.g.\ older or more ``difficult'' students placed together).
\item Confounder: with the same class duration, a different release time
means a different start time, so clock time (start and release) still
differs between treated and control. Time in school does not.
\end{itemize}


\noindent\textbf{(b) TWFE --- between schools \& within schools across
time, not-yet-treated control.}
\noindent\textbf{Level.} School.
 \noindent\textbf{Treatment.} school $s$ goes from some
shared release time with high school to zero shared release time.
Control: school $m$ still has some shared release time in both periods
(not-yet-treated).
Restriction: the shift of the 9th grade does not change between $t-1$
and $t$; the shift of high school can change.
\begin{equation}
  \hat\beta_{DD} = \big(\overline{V}_{s,t} - \overline{V}_{s,t-1}\big)
  - \big(\overline{V}_{m,t} - \overline{V}_{m,t-1}\big)
\end{equation}
$\overline{V}_{s,t}$ = average violence among 9th-grade students at school
$s$ in year $t$.

\begin{itemize}
\item Confounder:
 change in class-time attract more violent students for 9th grade.
\end{itemize}

\medskip
\noindent\textbf{(c) Dif-dif, school-shift level.}
\noindent\textbf{Level.} School $\times$ shift (morning/afternoon).
\begin{equation}
  \hat\beta_{DD} = \Big[\overline{V}^{\text{morning}}_{\text{shared}}
  - \overline{V}^{\text{morning}}_{\text{not shared}}\Big]
  - \Big[\overline{V}^{\text{afternoon}}_{\text{shared}}
  - \overline{V}^{\text{afternoon}}_{\text{not shared}}\Big]
\end{equation}
$\overline{V}^{\text{shift}}_{\text{shared}}$ = average violence among
9th-grade cohorts in that shift whose release time coincides with high
school's; $\overline{V}^{\text{shift}}_{\text{not shared}}$ = same, for
cohorts whose release time does not coincide.

\medskip
\noindent\textbf{(d) TWFE, student level.}
\noindent\textbf{Level.} Student, across their own grade progression (8th
$\to$ 9th). Student $m$ is a within-cohort control whose shared-release
status does not change from 8th to 9th grade.
\begin{equation}
  \hat\beta_{DD} = \big(Y_{s,\text{9th}} - Y_{s,\text{8th}}\big)
  - \big(Y_{m,\text{9th}} - Y_{m,\text{8th}}\big)
\end{equation}
$Y_{i,g}$ = violence outcome for student $i$ in grade $g$; $s$ = student
whose shared-release status changes moving from 8th to 9th grade; $m$ =
student whose shared-release status stays the same across both grades.



\subsection{Other effects}

\begin{itemize}
\item Other schools near. Affect in other schools
\item Problem: other schools can affect
\end{itemize}


